\ifx\danslelivre\undefined
\documentclass[10pt]{smfart}
\usepackage{ifthen}
\newboolean{english}
\newboolean{francais}
\setboolean{english}{false}
\setboolean{francais}{true}
\usepackage{palatino,euler}
\input{../fichiers-configuration-latex/paquets}
\selectlanguage{francais}
\input{../fichiers-configuration-latex/lemme-proposition-theoreme}
\input{../fichiers-configuration-latex/numerotation}
\input{../fichiers-configuration-latex/macros}
\input{../fichiers-configuration-latex/ucs_manquants}
\input{../.cv}

\externaldocument{../00-introduction/introduction}
\externaldocument{../01-quasi-excellents/excellent}
\externaldocument{../02-topologies/topologies}
\externaldocument{../03-approximation/approximation}
\externaldocument{../04-Cohen-Gabber/Cohen-Gabber}
\externaldocument{../05-algebrisation-partielle/algebrisation}
\externaldocument{../06-log-lissite/log}
\externaldocument{../08-modification-theorem/modification}
\externaldocument{../09-uniformisation-premiere-a-l/uniformisation-locale-premiere-a-l}
\externaldocument{../11-produits-orientes/produits-orientes}
\externaldocument{../12-descente-cohomologique-orientee/descente-orientee}
\externaldocument{../13-theoreme-de-finitude/finitude}
\externaldocument{../14-fonctions-dimension/fonctions-dimension}
\externaldocument{../15-Lefschetz-affine/Lefschetz-affine}
\externaldocument{../16-classes-de-Chern/gysin}
\externaldocument{../17-complexes-dualisants/dualite}
\externaldocument{../19-exemples-contre-exemples/contre-exemples}
\externaldocument{../20-rigidite/rigidite}
\externaldocument{../21-finitude-non-abelienne/finitude-non-abelienne}

\begin{document}

\begin{center}
VII. Démonstration du théorème d'uniformisation locale (faible)\\
Fabrice Orgogozo
\end{center}
\version

\else
\setboolean{english}{false}
\setboolean{francais}{true}
\selectlanguage{francais}
\chapter{Démonstration du théorème d'uniformisation
locale (faible)}
\addtocontents{toc}{par Fabrice Orgogozo}
\begin{center}
Fabrice Orgogozo
\end{center}
\fi

\section{Énoncé}

L'objet de cet exposé est de démontrer le théorème
\refext{II}{2-théorème uniformisation faible}
(voir aussi \refext{0}{0-uniformisation faible}),
dont nous rappelons l'énoncé ci-dessous :

\begin{théorème}\label{7-uniformisation}
Soient $X$ un schéma nœthérien quasi-excellent
et $Z$ un fermé rare de $X$.
Il existe une famille finie de morphismes 
$(X_i→X)_{i∈I}$, couvrante pour la topologie des
altérations et telle que pour tout $i∈I$ on ait :
\begin{enumerate}
\item le schéma $X_i$ est régulier et connexe ;
\item l'image inverse de $Z$ dans $X_i$ est le support
d'un diviseur à croisements normaux stricts.
\end{enumerate}
\end{théorème}

\section{Réductions : rappel des résultats antérieurs}\label{7-réductions}

\subsection{Réduction au cas local, normal de dimension finie}
Nous avons vu en \refext{II}{2-réduction cas hensélien}
qu'il suffit de démontrer le théorème lorsque
le schéma $X$ est local nœthérien normal hensélien excellent.
Faisons cette hypothèse supplémentaire. Un tel schéma est nécessairement de dimension finie,
que nous noterons ici $d$. De plus,
on a vu en \emph{loc. cit.} que si le théorème
est établi pour chaque schéma local nœthérien hensélien
excellent de dimension au plus $d$, il en est
de même pour les schémas nœthériens quasi-excellents
de dimension au plus $d$.

\subsection{Réduction au cas complet}
Il résulte de la proposition \refext{III}{3-reduc}
qu'il suffit de démontrer le théorème pour le schéma
local nœthérien complet $\chap{X}$, ce dernier
étant de même dimension que $X$ et également normal.

\subsection{Récurrence}
\label{7-récurrence}
Il résulte de ce qui précède que l'on peut supposer
le schéma $X$ local nœthérien complet normal
de dimension $d$ et le théorème connu pour chaque
schéma nœthérien quasi-excellent de dimension au plus $d-1$.
Lorsque $d=1$, le théorème est bien connu ; nous supposerons
dorénavant $d ≥ 2$.

\section{Fibration en courbes et application d'un théorème de A. J.
de Jong}


\subsection{}Soit $X=\Spec(A)$ un schéma local nœthérien complet normal
comme en \ref{7-récurrence} et $Z$ un fermé rare.
Quitte à remplacer $X$ (resp. $Z$) par un $X$-schéma fini
également local nœthérien normal excellent de dimension $d$
(resp. par son image inverse), on peut supposer
d'après \refext{V}{5-algébrisation partielle mixte raffinée géométrique},
qu'il existe un schéma local nœthérien régulier $S$ de dimension $d-1$,
un $S$-schéma de type fini dominant $X′$ intègre et affine, un point fermé $x′$ de la
fibre spéciale de $f:X′ → S$, et un fermé rare $Z′$ de $X′$
satisfaisant les conditions suivantes :

— le morphisme $c$ induit un isomorphisme $X ⥲ \Spec(\chap{𝒪_{X′,x′}})$ ;

— l'image inverse $c^{-1}(Z′)$ de $Z′$ coïncide avec $Z$.

\begin{center}
\begin{tikzpicture}[auto]
\matrix(diag)[matrix of math nodes]
{|(X)| X & |(Xprime)| X′ \\
   & |(S)| S  \\
};
 \draw[->] (X) -- node{$c$}  (Xprime);
 \draw[->] (Xprime) -- node{$f$} (S);
\end{tikzpicture}
\end{center}

\subsection{}Supposons l'existence d'une famille $(X′_i → X′)$
couvrante pour la topologie des altérations (\refext{II}{2-définition altL})
telle que chaque $X_i′$ soit régulier et chaque image inverse $Z_i ′$
de $Z′$ dans $X_i′$ soit le support d'un diviseur à croisements
normaux. Il résulte de \refext{II}{2-stabilité par cb}
que la famille $(X_i → X)$ obtenue par changement de base
(plat) $X → X′$ est également $\alt$-couvrante.
D'autre part, l'hypothèse d'excellence faite sur les schémas
garantit que le morphisme de complétion $c$ est \emph{régulier}
(\refext{I}{1def-excellent}).
La régularité d'un morphisme étant stable par changement de base
localement de type fini (\ega{IV₂}{6.8.3}),
et préservant la régularité des schémas (\ega{IV₂}{6.5.2 (ii)})
il en résulte que chaque $X_i$ est régulier. De même, l'image
inverse $Z_i$ de $Z_i′$ dans $X_i$ — qui coïncide
avec l'image inverse de $Z$ dans $X_i$ par le morphisme
évident — est le support d'un diviseur à croisements normaux
pour chaque indice $i$. 

\subsection{}Quitte à remplacer $X$ par $X′$, ce
qui est licite d'après ce qui précède, nous pouvons
supposer le schéma $X$ intègre de dimension $d$,
équipé d'un morphisme dominant de type fini $f:X → S$ où $S$
est local nœthérien régulier de dimension $d-1$.
Quitte à compactifier $f$, on peut le supposer \emph{propre} ;
quitte à éclater, on peut supposer que le fermé $Z$ est un
\emph{diviseur}.

\subsection{}
Nous sommes dans les conditions d'application
du théorème \cite[2.4]{Families@de_Jong},
d'après lequel, quitte à altérer $S$ et $X$, on peut supposer les faits suivants :

— le morphisme $f$ est une courbe nodale ;

— le diviseur $Z$ est contenu dans la réunion
d'un diviseur $D$ étale sur $S$, contenu dans le lieu
lisse de $f$, et de l'image inverse $f^{-1}(T)$ d'un
fermé rare $T$ de $S$.

\section{Résolution des singularités}
\label{7-utilisation log-résolution}

\subsection{Résolution des singularités de la base}
Les altérations précédentes conduisent à une situation
où le schéma $X$ n'est plus local (ni même affine) et
le schéma $S$ n'est plus nécessairement régulier.
Il est cependant excellent de dimension $d-1$
donc justiciable de l'hypothèse de récurrence \ref{7-récurrence}.
Ainsi, on peut supposer que la paire $(S,T)$ est régulière,
c'est-à-dire que le schéma $S$ est régulier et que $T$
est un diviseur à croisements normaux. En effet,
il en est ainsi localement pour la topologie des altérations.

\subsection{}
D'après \refext{VI}{6-nodal implique log lisse},
la paire $(X,D ∪ f^{-1}(T))$ est log-régulière
au sens de \refext{VI}{6-1.2}.
Observons qu'un diviseur contenu dans un
diviseur à croisements normaux est lui-même à croisements normaux.
Ceci nous permet de supposer d'agrandir $Z$ si
nécessaire et auquel on applique alors le théorème suivant
de Katô K. (\cite[10.3, 10.4]{Toric@Kato}), complété par W. Nizioł
(\cite[5.7]{Toric@Niziol}). (Voir aussi \cite[§5.6]{HT@GR} et \refext{VIII}{8-resolution of log schemes1}.)

\begin{théorème}
Soit $(X,Z)$ une paire log-régulière, où $X$ est un schéma nœthérien.
Il existe un schéma noethérien régulier $Y$, un diviseur à
croisements normaux $D⊂Y$ et un morphisme projectif birationnel $π:Y→ X$
tel que l'image inverse $π^{−1}(Z)$ soit contenue dans $D$.
\end{théorème}

\ifx\danslelivre\undefined
\bibliography{../fichiers-configuration-latex/bibliographie}
\bibliographystyle{../fichiers-configuration-latex/nomcomplet-date}
\end{document}
\fi
